\section{The finite path-system argument}\label{sec:finite}

It remains to treat maximum cliques of order at most eight.  The clique
bound in Theorem~\ref{clq:bound} makes this a finite problem: if
$\lfloor k/2\rfloor\le t\le8$, then $k\le2t+1\le17$.
We shall rule out the optimal path systems of Section~\ref{sec:paths}
by a finite list of deductions from their required edges.  The objects
enumerated are weighted path systems; the ambient graphs are not
enumerated.

\begin{proposition}[Finite verification]\label{finite:verified}
Let $k,t$ be integers such that
\begin{equation}\label{finite:domain}
 5\le k\le17,\qquad
 \max\{3,\lfloor k/2\rfloor\}\le t\le\min\{8,k\}.
\end{equation}
There is no finite simple graph $G$ satisfying all of the following:
$G$ has no cycle on exactly $k+1$ vertices, its clique number is $t$,
$\alpha(G)\le k$, and
\[
 |N_G[I]|\ge k|I|+1
 \quad\text{for every nonempty independent set }I\subseteq V(G).
\]
More precisely, every optimal path system on a maximum clique, chosen
as in Section~\ref{sec:paths}, gives a contradiction under these
hypotheses.
\end{proposition}

The proof will be completed after we establish the deductions used by
the computation and give its verified results.  Throughout this section,
suppose that such a graph exists.  Fix a maximum clique $Q$ of order $t$,
put $h=k+1-t$, and choose a path system $\mathcal P$ of maximum total
amount $L<h$, with the number $e$ of its paths as small as possible.
An amount is the number of internal vertices of the corresponding path.
In particular,
\begin{equation}\label{finite:budget}
 0\le L\le k-t,\qquad
 S=Q\cup\bigcup_{R\in\mathcal P}\bigl(V(R)\setminus Q\bigr),
 \qquad |S|=t+L\le k.
\end{equation}
The expansion hypothesis gives $\delta(G)\ge k$.  We also retain the
forbidden amount interval from Lemma~\ref{path:interval}, the separation
rules of Lemma~\ref{path:separation}, and the size test of
Lemma~\ref{size:test}.

\subsection{Patterns and their complete enumeration}

The paths of $\mathcal P$ represent the edges of a linear forest on
$Q$.  Include every unused vertex of $Q$ as a singleton component of
this forest.  Orient a component temporarily and list its edge amounts
in order, obtaining a tuple $p=(q_1,\ldots,q_r)$ of positive integers.
A singleton has the empty tuple $()$.  Reversing the component reverses
its tuple, so we choose the lexicographically smaller of these two
tuples.  Sort the resulting component tuples in nondecreasing
lexicographic order, retaining repeated components.  Their list
\[
 P=(p_1,\ldots,p_c)
\]
is the \emph{pattern} of the path system.  For a tuple $p$, write
$\ell(p)$ for its number of entries and $\sigma(p)$ for their sum,
with both quantities zero when $p=()$.  Thus
\begin{equation}\label{finite:pattern-parameters}
 \sum_{i=1}^c\bigl(\ell(p_i)+1\bigr)=t,\qquad
 \sum_{i=1}^c\sigma(p_i)=L\le k-t,\qquad
 \sum_{i=1}^c\ell(p_i)=e.
\end{equation}

\begin{lemma}[Pattern coverage]\label{finite:coverage}
For nonnegative integers $t,B$, the nondecreasing lists of
reversal-normalized positive tuples satisfying
\[
 \sum_{p\in P}(\ell(p)+1)=t,
 \qquad \sum_{p\in P}\sigma(p)\le B
\]
represent exactly once the isomorphism classes of linear forests on
$t$ vertices with positive integer edge amounts of total at most $B$.
Here empty tuples are allowed as singleton components.

All these lists are generated by the following recursion.  Choose a
first tuple $p$ with $\ell(p)\le t-1$ and $\sigma(p)\le B$, subject
to reversal normalization and any prescribed lower bound on $p$.
Continue with $t-\ell(p)-1$ vertices, budget $B-\sigma(p)$, and
lower bound $p$ on the next tuple.  At zero remaining vertices, emit
the empty list for every remaining nonnegative budget.
\end{lemma}

\begin{proof}
An isomorphism of two nontrivial paths either preserves or reverses
their endpoint order.  Therefore their weighted isomorphism classes
agree exactly when their amount tuples agree up to reversal.  A forest
isomorphism permutes components and acts in this way within each one.
Reversal normalization followed by sorting consequently gives a
complete invariant, with multiplicities preserved.

Every candidate first tuple is obtained by listing positive integer
compositions of each amount from zero to $B$, restricting the number
of entries to at most $t-1$.  The amount-zero candidate is just the
empty tuple.  For example, one may start with the empty tuple and
repeatedly append a positive integer no larger than the remaining
budget.  Each tuple is reached once.  Normalization and the lower bound
are tested when selecting a component; a rejected prefix must still
be extended, since a longer tuple can satisfy these conditions even
when its prefix does not.

Now induct on $t$.  For $t=0$ the empty list is the unique pattern,
regardless of unused budget.  For $t>0$, the first tuple consumes
$\ell(p)+1\ge1$ vertices and $\sigma(p)$ amount.  The inductive
hypothesis supplies every allowed tail exactly once, while its lower
bound $p$ enforces nondecreasing order.  Conversely, every permitted
list has exactly this first tuple and one of these tails.  This proves
both coverage and uniqueness.
\end{proof}

In particular, a system containing no paths is represented by $t$
empty component tuples.  The outer empty list instead represents a
forest on zero vertices.  Unused amount budget is allowed: optimality
does not require $L=k-t$.  Enumerating some patterns that cannot occur
as optimal systems causes no difficulty; Lemma~\ref{finite:coverage}
ensures that every system that could occur is present.

There are exactly $42$ parameter pairs in \eqref{finite:domain},
and their pattern lists contain $3099$ instances in total.  A pattern
appearing for two different parameter pairs is counted once for each
pair.  Table~\ref{finite:counts} gives the counts.  For completeness,
these numbers have a separate elementary enumeration.  If $a(r,m)$
denotes the number of chain types with $r\ge1$ edges and amount
$m\ge r$, reversal pairs all positive compositions except the
palindromes.  Hence
\begin{equation}\label{finite:chain-count}
 a(r,m)=\frac12\left(\binom{m-1}{r-1}+b(r,m)\right),
\end{equation}
where
\[
 \begin{aligned}
 b(2s,m)&=
 \begin{cases}
  \binom{m/2-1}{s-1},&m\text{ even},\\
  0,&m\text{ odd},
 \end{cases}\\
 b(2s+1,m)&=\binom{\lfloor(m-1)/2\rfloor}{s}.
 \end{aligned}
\]
The even formula follows by pairing equal entries; the odd formula
also allows a positive central entry and sums over the paired entries.
The case $s=0$ in the odd formula gives one type with a single edge.
Choosing a multiset of chain types and then singleton components shows
that the count for $(t,B)$ is
\begin{equation}\label{finite:count-series}
 \sum_{L=0}^{B}[x^t y^L]\,
 \frac{1}{1-x}
 \prod_{r\ge1}\prod_{m\ge r}
       (1-x^{r+1}y^m)^{-a(r,m)}.
\end{equation}
Only factors with $r+1\le t$ and $m\le B$ contribute.  Thus the
table can be checked by finite integer arithmetic independently of
the deductions that exclude the patterns.

\begin{table}[htbp]
\centering
\begin{tabular}{ccl}
$k$ & admissible $t$ & number of patterns, summed over $t$\\
\hline
5 & $3,4,5$ & $4+2+1=7$\\
6 & $3,4,5,6$ & $6+5+2+1=14$\\
7 & $3,4,5,6,7$ & $9+9+5+2+1=26$\\
8 & $4,5,6,7,8$ & $16+10+5+2+1=34$\\
9 & $4,5,6,7,8$ & $25+20+11+5+2=63$\\
10 & $5,6,7,8$ & $35+24+11+5=75$\\
11 & $5,6,7,8$ & $60+46+25+11=142$\\
12 & $6,7,8$ & $87+51+26=164$\\
13 & $6,7,8$ & $152+104+55=311$\\
14 & $7,8$ & $197+118=315$\\
15 & $7,8$ & $364+237=601$\\
16 & $8$ & $468$\\
17 & $8$ & $879$
\end{tabular}
\caption{All admissible parameter pairs and their pattern counts,
with amount budget $B=k-t$.  These are counts of weighted forests,
not of ambient graphs.}
\label{finite:counts}
\end{table}

\subsection{Required edges and path labels}

We now give each pattern a concrete graph on $S$, to which the
subsequent deductions can be applied.  Process its components in
their chosen order and traverse each one in its chosen orientation.
Label its clique vertices and all intervening path interiors
consecutively, beginning with label zero for the first component and
continuing with the next unused label at each new component.
Thus $S$ is identified with $\{0,\ldots,t+L-1\}$.

Let $\mathcal E(P)$ be the ordered list of pairs $(a,b)$ of clique
labels corresponding to the original paths, in this traversal order.
An edge of amount $q$ has $b=a+q+1$, and its expanded path is
$a,a+1,\ldots,b$.  Let $J$ be the graph whose edges are all clique
edges on $Q$ and all consecutive steps of these expanded paths.
We call these the \emph{required edges}: every one is present in
$G[S]$, while a pair not joined in $J$ may or may not be an edge
of $G[S]$.

For a label $x\in S$, define the set of endpoint choices
\begin{equation}\label{finite:tips}
 T(x)=
 \begin{cases}
  \{x\},&x\in Q,\\
  \{a,b\},&a<x<b\text{ for }(a,b)\in\mathcal E(P).
 \end{cases}
\end{equation}
For $x\in Q$, the sole choice represents the segment of length zero.
For an interior vertex, the two choices represent the two segments
along its original path to its clique endpoints.

\begin{lemma}[Label construction]\label{finite:labels}
The construction labels every vertex of $S$ exactly once.  It has
exactly $t$ clique labels and $L$ interior labels, and
\[
 |\mathcal E(P)|=e,\qquad
 \sum_{(a,b)\in\mathcal E(P)}(b-a-1)=L.
\]
Every interior label belongs to exactly one open interval $(a,b)$
from $\mathcal E(P)$, so \eqref{finite:tips} is well defined.
The list $\mathcal E(P)$ is increasing, with $a<b$ in every pair.
For any subset of its edges, a connecting route from a clique label
$z$ to a larger clique label $w$ follows increasing edges in list order.
\end{lemma}

\begin{proof}
A component tuple $p$ occupies exactly
$\ell(p)+1+\sigma(p)$ consecutive labels.  Its successive expanded
paths share their common clique endpoint and have disjoint interiors.
Different components occupy disjoint intervals of labels.  Summing
these facts gives the asserted counts and the unique interval for
each interior vertex, including when some or all components are
singletons.

Within a component, the clique labels strictly increase along the
path; all labels of an earlier component precede those of a later
one.  A subset of the old edges is a union of subpaths of these
components.  Its unique route between connected labels $z<w$
therefore traverses increasing labels, and these edges appear in
the same order in $\mathcal E(P)$.
\end{proof}

Relabeling and reversing components of an actual system produces
exactly this required-edge representation.  Additional ambient
edges are retained as possibilities throughout the argument.  We
next use the optimality of $\mathcal P$ and the absence of a
$C_{k+1}$ to deduce which paths outside $S$ are impossible.
